%% AISSH-26 workshop expansion manuscript: Springer Nature sn-jnl format.
%% V2 -- NARRATIVE RESTRUCTURE of aissh_springer_manuscript_12pg.tex. Same
%% numeric results, tables, and citations throughout; no new experiments, no
%% strengthened claims. What changed is framing and order: the paper now leads
%% with one central story ("we built a deployable, age-adaptive pediatric
%% conversational system, and in doing so uncovered two scientific findings"),
%% with guard-bear's safety-transfer failure as the strongest contribution,
%% consumer-hardware deployability second, and the LoRA configuration-ordering
%% crossover reframed as an unexpected discovery made during development
%% (existence established, mechanism only partially resolved) rather than the
%% paper's primary motivation. Results are reordered (response model ->
%% guard-bear -> deployment synthesis -> multi-seed statistical validation ->
%% mechanism sweeps -> cross-architecture), Related Work subsections reordered
%% to match, and the previously-scattered reproducibility narrative (dataset
%% revision, timestamp audit, library-version check, run-to-run variance) is
%% consolidated into one Methods subsection instead of being repeated in
%% Results, Discussion, and Conclusion. Claims softened where the evidence is
%% an automated proxy, not ground truth (e.g. "age-appropriate" ->
%% "age-adapted" / "automated proxies for"); mechanism language kept
%% conservative ("suggests", "consistent with", "partially resolved").
%% The original 12-page submission draft is preserved untouched at
%% aissh_springer_manuscript_12pg.tex -- this file is a separate version.
%%
%% TODO before submission:
%%   - Fill in real institutional emails for all four authors (FIXME markers below).
%%   - Confirm venue/copyright/funding boilerplate in the Declarations section.
%%   - Confirm whether AISSH-26 permits the repository link in the appendix.
%%   - Recompile and check final page count once on the actual submission
%%     toolchain/template version (this copy was compiled against a GitHub
%%     mirror of sn-jnl.cls, not the official v3.1 Dec-2024 package).

\documentclass[pdflatex,sn-mathphys-num]{sn-jnl}% Math and Physical Sciences Numbered Reference Style

\usepackage{graphicx}%
\usepackage{multirow}%
\usepackage{amsmath,amssymb,amsfonts}%
\usepackage{amsthm}%
\usepackage{mathrsfs}%
\usepackage{xcolor}%
\usepackage{textcomp}%
\usepackage{manyfoot}%
\usepackage{booktabs}%
\usepackage{enumitem}%

\raggedbottom
\renewcommand{\arraystretch}{0.85}
\setlist{nosep}

\begin{document}

\title[Safety-Gated, Age-Adaptive LoRA Fine-Tuning for a Virtual Pediatric Hospital Guide]{Safety-Gated, Age-Adaptive LoRA Fine-Tuning for a Virtual Pediatric Hospital Guide}

\author*[3]{\fnm{Anthony} \sur{Barros}}\email{FIXME-institutional-email@uwo.ca}

\author[1,2]{\fnm{Sandrine} \sur{de Ribaupierre}}\email{FIXME-institutional-email@uwo.ca}

\author[2,3]{\fnm{Roy} \sur{Eagleson}}\email{FIXME-institutional-email@uwo.ca}

\author[3]{\fnm{Kelsey} \sur{Kloosterman}}\email{FIXME-institutional-email@uwo.ca}

\affil*[1]{\orgdiv{Clinical Neurological Sciences}, \orgname{Western University}, \orgaddress{\city{London}, \state{Ontario}, \country{Canada}}}

\affil[2]{\orgdiv{Centre for Brain and Mind}, \orgname{Western University}, \orgaddress{\city{London}, \state{Ontario}, \country{Canada}}}

\affil[3]{\orgdiv{Department of Artificial Intelligence and Software Engineering, Faculty of Engineering}, \orgname{Western University}, \orgaddress{\city{London}, \state{Ontario}, \country{Canada}}}

\abstract{Pediatric patients navigating hospital environments need age-appropriate, institutionally grounded information that general-purpose AI systems are not designed to provide, and any system that accepts free-text input from child users must also defend against inputs it was never designed to handle. We present \textit{Dr.\ Beary Good}, an age-targeted conversational guide deployed within a digital recreation of Victoria Hospital (London, Ontario), paired with \textit{guard-bear}, an upstream input-safety classifier gating access to it -- together forming a complete pipeline that runs entirely on consumer hardware (an Apple Mac Mini M4, 16~GB) with no cloud dependency. The first major finding concerns safety: general-purpose prompt-injection classifiers do not describe the input distribution a pediatric conversational system actually receives. The unmodified base classifier we build on, \texttt{Prompt-Guard-86M}~\cite{promptguard2024}, performs \textit{worse than random} (ROC-AUC 0.466) on a held-out set of pediatric-context queries, saturating by flagging nearly everything as unsafe (0.505 precision despite 0.959 recall). Fine-tuning \textit{guard-bear} on 4{,}703 pediatric-specific labeled examples (14 subcategories) resolves this to near-perfect discrimination (ROC-AUC 0.999, 0.987 recall, 0.990 precision) -- evidence that safety mechanisms validated for a developer-facing threat model cannot be assumed to transfer to a different deployment population without dedicated evaluation. The second finding concerns the response model. We fine-tune Llama 3.2 Instruct (1B, 3B) with LoRA at two configurations -- Standard ($r=8$, 16 layers) and Fast ($r=4$, 8 layers) -- across both age groups (5--11, 12--18) on 1123 LIMA-inspired~\cite{zhou2023lima} synthetic examples. Both configurations substantially outperform zero-shot base models on automated readability proxies (56--66\% FK~$\leq$~7.0 pass rate vs.\ 12\%) and inter-role lexical separation (0.92--0.96 vs.\ 0.64--0.67 classifier accuracy), and Fast-1B meets the 1.0-second real-time VR latency target. During development, an initial single-seed comparison suggested an unexpected \textbf{configuration-ordering crossover} (Standard better on 3B, Fast better on 1B); because a routine reproducibility check found this pipeline's run-to-run variance can exceed that gap, we validated the observation with a dedicated 110-run multi-seed campaign, confirming it is statistically real on both sides -- though a length-adjustment check shows the 3B-side gap is substantially a response-length effect, while the 1B-side gap survives as a length-independent difference. Follow-up rank and layer sweeps partially, but not fully, attribute the crossover's mechanism to LoRA capacity and depth. Together, these results show that a deployable, age-adaptive pediatric conversational pipeline is achievable on consumer hardware, and that its safety layer -- not its response quality -- is where an unvalidated general-purpose component would have posed the larger practical risk.}

\keywords{parameter-efficient fine-tuning, low-rank adaptation, pediatric clinical NLP, virtual reality, on-device inference, small language models, LLM input safety classification}

\maketitle

\section{Introduction}\label{sec1}

Large language models (LLMs) enable domain-specific conversational systems, but resource-constrained clinical deployment raises constraints rarely addressed together: limited memory favors smaller models, domain adaptation requires fine-tuning on specialized data, real-time interaction imposes strict latency bounds, and free-text input from a vulnerable population must be defended against inputs the system was never designed to handle.

This paper presents \textit{Dr.\ Beary Good}, an age-targeted conversational guide embedded in a virtual recreation of Victoria Hospital (London, Ontario), providing age-appropriate hospital information to children aged 5--11 and adolescents aged 12--18 within a VR environment, together with \textit{guard-bear}, an upstream input-safety classifier gating access to it. All training and inference runs on an Apple Mac Mini M4 (16~GB): a complete age-adaptive pipeline -- response generation and input safety together -- deployable on consumer hardware without cloud infrastructure.

A conversational system that accepts open-ended text from children needs a defense layer distinct from the response model itself. General-purpose prompt-injection classifiers such as \texttt{Prompt-Guard-86M}~\cite{promptguard2024} target a developer-facing threat model that does not describe a pediatric-facing system's actual inputs: short, first-person, emotionally direct legitimate queries on one side, out-of-scope curiosity, gibberish, and adult/parental content on the other. We are not aware of prior work evaluating an input-safety classifier calibrated to this population (Section~\ref{sec:relatedwork}); \textit{guard-bear} closes this gap, and finds, as our central safety result, that the unmodified base classifier performs worse than random here, a failure fine-tuning fully resolves.

For the response model, we fine-tune Llama 3.2 Instruct models~\cite{dubey2024llama3} at 1B and 3B scales using LoRA~\cite{hu2022lora}, standard practice for resource-constrained adaptation alongside QLoRA~\cite{dettmers2023qlora}. Training data follows the LIMA hypothesis~\cite{zhou2023lima}: $\sim$560 examples per adapter (1123 total across six categories) grounded in Victoria Hospital's public materials. We compare Standard ($r=8$, 16 layers) and Fast ($r=4$, 8 layers) adapters across both sizes and age groups (eight runs total). Pediatric deployment introduces demands distinct from adult-facing clinical LLMs and passive VR distraction tools alike -- developmental vocabulary, emotionally supportive register, hard safety constraints, responsiveness to patient queries -- that neither generic fine-tuning nor scripted content delivers; we are not aware of prior work fine-tuning age-stratified LoRA adapters for this setting (Section~\ref{sec:relatedwork}).

During development of this response model, an initial single-seed ablation suggested an unexpected \textbf{configuration-ordering crossover}: Standard yields better readability and style separation on 3B, while Fast outperforms on 1B -- the opposite of what existing rank-selection guidance, developed only at the 7B--70B scale, would predict. A routine reproducibility check found run-to-run training variance exceeding the crossover gap itself, so we validated it directly with a dedicated 110-run multi-seed campaign (Section~\ref{sec:seedcampaign}), confirming it is statistically real on both sides. A subsequent rank sweep (Section~\ref{sec:ranksweep}) sought to attribute it to rank as the operative variable, replicating cleanly for one architecture (SmolLM2) but only weakly for others. Only Fast-1B meets the 1.0-second real-time VR latency target among Llama configurations for the 5--11 age group; several cross-architecture configurations also clear this target (Section~\ref{sec:crossarch}).

\medskip\noindent The contributions of this paper are:
\begin{itemize}
    \item \textit{guard-bear}: a fine-tuned input-safety classifier for the pediatric clinical threat model, showing its unmodified general-purpose base model performs \textit{worse than random} (ROC-AUC 0.466) on this population -- resolved by fine-tuning to near-perfect discrimination (ROC-AUC 0.999) -- evidence that general-purpose safety classifiers cannot be assumed to transfer to a new deployment population without dedicated evaluation.
    \item A complete age-adaptive pipeline -- response generation gated by an upstream safety classifier -- trained and evaluated end-to-end on consumer hardware with no cloud dependency, showing pediatric conversational deployment does not require server-scale infrastructure.
    \item An unexpected LoRA configuration-ordering crossover between Llama 3.2 1B and 3B, discovered during development and confirmed statistically real on both sides by a dedicated 110-run multi-seed campaign (1B: $t=6.40$, $p=6.0\times10^{-8}$, $n=25$/config; 3B: $n=30$/config, corrected $p=0.0015$ via a combination test after an interim sample-size look) after an initial single-seed observation proved non-reproducible at face value. A length-adjustment check further shows the 1B-side gap is a length-independent register effect, while the 3B-side gap is substantially attributable to response length.
    \item A rank sweep (32 runs) and layer sweep (12 runs) investigating the crossover's mechanism: the layer-depth contribution replicates reasonably well across an independent two-seed run, while the rank-regime story replicates cleanly only for SmolLM2 and is much weaker for Llama and Qwen at two-seed power -- evidence that small-model LoRA rank selection below 7B is more architecture-dependent than existing guidance assumes, detailed in Section~\ref{sec:ranksweep}.
    \item Cross-architecture evidence on Qwen~2, SmolLM2, and Qwen 2.5 (0.5B--3B), all re-validated on one consistent dataset vintage, that Fast-vs-Standard directional patterns generalize beyond Llama even where the underlying rank-regime mechanism does not cleanly resolve.
\end{itemize}

\section{Related Work}\label{sec:relatedwork}

\subsection{Input-Safety Classifiers for Deployed LLM Systems}

General-purpose prompt-injection classifiers such as \texttt{Prompt-Guard-86M}~\cite{promptguard2024} target developer-facing override patterns. The closest child-safety effort, \textit{KidRails}~\cite{kidrails2025}, fully fine-tunes an 8B response model on curated safe/unsafe pairs, without a separate classifier or baseline comparison. \textit{guard-bear} instead evaluates a dedicated input classifier against its base model on a labeled, subcategorized pediatric threat taxonomy -- to our knowledge, the first such evaluation for this population.

\subsection{LLMs for Pediatric and Clinical Communication}

Prompted, not fine-tuned, LLMs have been evaluated as pediatric communication aids: Rao et al.~\cite{rao2025pediatric} found off-the-shelf models can produce age-appropriate surgical explanations but called for further fine-tuning; a Norwegian evaluation similarly found generated child-directed language skewed too advanced for the target age~\cite{hassan2025childconv}. \textit{KidLM}~\cite{nayeem2024kidlm} adapts a masked-language-model pretraining objective to a children's corpus but targets general language modeling, not age-stratified clinical generation. No prior work, to our knowledge, fine-tunes age-stratified LoRA adapters for pediatric hospital communication. VR-based pediatric interventions reduce procedural anxiety~\cite{ryu2021vr}, but existing deployments are largely passive, scripted content; \textit{Dr.\ Beary Good} integrates fine-tuned, age-targeted adapters into an interactive VR environment instead.

\subsection{Parameter-Efficient Fine-Tuning at Small Scale}

LoRA~\cite{hu2022lora} and QLoRA~\cite{dettmers2023qlora} are standard for resource-constrained fine-tuning, but rank-selection guidance -- including Biderman et al.'s~\cite{biderman2024lora} empirical study, the most thorough to date -- is developed and validated at the 7B--70B scale; rank-ordering behavior below 7B is uncharacterized by prior work. Our rank and layer sweeps (Section~\ref{sec:ranksweep}) address this gap directly.

\section{Results}\label{sec2}

This section first establishes the response model and \textit{guard-bear} as deployable components (Sections~\ref{sec:responsemodel}--\ref{sec:deployment}), then presents the statistical validation and mechanistic investigation of an unexpected LoRA effect discovered during development (Sections~\ref{sec:seedcampaign}--\ref{sec:crossarch}). Methods and dataset detail: Section~\ref{sec:methodology}.

\subsection{Response Model: Perplexity, Readability, Latency, and Style Separation}\label{sec:responsemodel}

Table~\ref{tab:res2} reports a representative single run (seed 42) of the four fine-tuned configurations and both base models, current dataset vintage (Section~\ref{subsec:responsedataset}). All four fine-tuned adapters substantially outperform base models on FK~$\leq$~7.0 for the 5--11 group (56--66\% vs.\ 12\%), driven by shorter responses (74--99 vs.\ 178--185 avg words), and on inter-role style separation (0.92--0.96 vs.\ 0.64--0.67; TF-IDF+logistic-regression, 5-fold CV, chance~=~0.50) -- though this classifier does not discriminate Standard from Fast, consistent with, not independent confirmation of, the readability pattern below. \textbf{Fast-1B is the only Llama configuration meeting the 1.0s VR latency target} (0.97s avg, 62\% under target; several cross-architecture configurations, Table~\ref{tab:crossarch}, also clear it); all 3B configurations exceed 2.2s. Post-hoc perplexity does not track FK-based readability cleanly: Fast has lower perplexity than Standard at \emph{both} sizes (1B: 30.18 vs.\ 30.46; 3B: 22.91 vs.\ 24.88).

On this seed, Fast-1B beats Standard-1B on FK pass rate (62\% vs.\ 56\%), while Standard-3B and Fast-3B tie (66\% vs.\ 66\%) -- an apparent \textbf{configuration-ordering crossover} between sizes; Section~\ref{sec:seedcampaign} shows this single-run picture undersells the true (smaller, but real) 3B-side gap.

\begin{table*}[!h]
  \caption{Readability (50 age-matched held-out prompts/group) and latency/throughput (sequential, Mac Mini M4), single run (seed 42), current dataset vintage.}
  \label{tab:res2}
  \scriptsize
  \setlength{\tabcolsep}{3pt}
  \begin{tabular}{@{}llcccccccc@{}}
    \toprule
    Config. & Age & FK & SMOG & G.Fog & C-L & TTR & FK$\leq$7.0 & Avg Lat.(s) & $<$1.0s \\
    \midrule
    Standard-1B & 5--11  & 7.1  & 9.2  & 8.7  & 7.1  & 0.756 & 56\% & 1.12 & 26\% \\
    Standard-1B & 12--18 & 10.3 & 12.7 & 13.0 & 10.7 & 0.734 & ---  & 1.36 & 2\%  \\
    Standard-3B & 5--11  & 6.6  & 8.9  & 8.1  & 7.1  & 0.781 & 66\% & 2.22 & 0\%  \\
    Standard-3B & 12--18 & 10.4 & 12.8 & 12.9 & 10.6 & 0.733 & ---  & 3.08 & 0\%  \\
    Fast-1B     & 5--11  & 6.5  & 9.0  & 8.0  & 6.8  & 0.752 & 62\% & 0.97 & 62\% \\
    Fast-1B     & 12--18 & 9.7  & 12.2 & 12.1 & 10.0 & 0.725 & ---  & 1.11 & 38\% \\
    Fast-3B     & 5--11  & 6.5  & 9.0  & 8.1  & 7.1  & 0.724 & 66\% & 2.70 & 0\%  \\
    Fast-3B     & 12--18 & 10.0 & 12.5 & 12.5 & 10.4 & 0.663 & ---  & 3.63 & 0\%  \\
    Base-1B     & 5--11  & 9.5  & 11.9 & 12.0 & 9.6  & 0.582 & 12\% & 1.90 & 16\% \\
    Base-1B     & 12--18 & 10.4 & 12.9 & 13.1 & 11.6 & 0.553 & ---  & 2.23 & 2\%  \\
    Base-3B     & 5--11  & 9.8  & 12.4 & 12.5 & 10.2 & 0.609 & 12\% & 4.54 & 8\%  \\
    Base-3B     & 12--18 & 11.7 & 13.6 & 14.7 & 11.4 & 0.551 & ---  & 5.40 & 2\%  \\
    \bottomrule
  \end{tabular}
\end{table*}

\subsection{guard-bear: Safety Classification Results}\label{sec:guardresults}

The clearest evidence that general-purpose safety tooling cannot be assumed to transfer to a new population comes from comparing \textit{guard-bear} against unmodified \texttt{Prompt-Guard-86M} (default threshold 0.50) on the 613-example held-out test set (300 safe, 313 unsafe; Table~\ref{tab:guardheadline}). \textbf{The base model is worse than random on this population}: ROC-AUC 0.466 is below the 0.50 chance baseline, since its imperative-override feature space does not describe short, first-person pediatric queries or the gibberish/out-of-scope/adult content that should also be blocked. Its apparent 0.9585 recall is an artifact of saturation, flagging 294 of 300 safe examples as unsafe too (precision 0.505); a sub-0.5 AUC is stronger than saturation alone, meaning the base model ranks true-unsafe examples below true-safe examples across the full threshold range, not just that 0.50 is miscalibrated. \textbf{Fine-tuning resolves the discrimination failure}, not merely the threshold: ROC-AUC rises to 0.999, false negatives fall from 13 to 4, false positives from 294 to 3, meeting all four deployment targets (recall $\geq$0.97, precision $\geq$0.85, F1 $\geq$0.90, accuracy $\geq$0.92) with margin. The largest subcategory gain is on \texttt{gibberish\_malformed} (48\%~$\to$~100\% recall); both models reach 100\% on the canonical adversarial categories the base model was designed to catch.

\begin{table}[!h]
  \centering
  \caption{guard-bear vs.\ base model, held-out test set (613 examples: 300 safe, 313 unsafe).}
  \label{tab:guardheadline}
  \small
  \begin{tabular}{@{}lccc@{}}
    \toprule
    Metric & Base (t=0.50) & guard-bear (t=0.08) & $\Delta$ \\
    \midrule
    Unsafe recall    & 0.9585 & \textbf{0.9872} & $+$0.029 \\
    Unsafe precision & 0.5051 & \textbf{0.9904} & $+$0.485 \\
    Unsafe F1        & 0.6615 & \textbf{0.9888} & $+$0.327 \\
    Accuracy         & 0.4992 & \textbf{0.9886} & $+$0.489 \\
    ROC-AUC          & 0.4660 & \textbf{0.9993} & $+$0.533 \\
    \bottomrule
  \end{tabular}
  \qquad
  \begin{tabular}{@{}lcc@{}}
    \toprule
    Confusion & Base & guard-bear \\
    \midrule
    TN & 6   & 297 \\
    FP & 294 & 3   \\
    FN & 13  & 4   \\
    TP & 300 & 309 \\
    \bottomrule
  \end{tabular}
\end{table}

\label{sec:deployment}Table~\ref{tab:res2}'s latency figures and \textit{guard-bear}'s footprint (86M parameters, full fine-tune) together show a complete pipeline -- safety gating then age-routed response generation -- practical on one consumer machine; Fast-1B alone clears the 1.0s VR target (0.97s avg), as do several cross-architecture configurations (Table~\ref{tab:crossarch}). End-to-end latency with \textit{guard-bear}'s pass included has not yet been measured (Limitations); the two components have only been benchmarked separately.

\subsection{Statistical Validation of the Crossover: A Multi-Seed Campaign}\label{sec:seedcampaign}

The single-run pattern above (Section~\ref{sec:responsemodel}) was an unexpected finding uncovered during development, not a designed comparison, so it needed the same reproducibility scrutiny as any surprising result before being treated as real. The ablation and the rank/layer sweeps in Section~\ref{sec:ranksweep} are single- or two-seed comparisons; Section~\ref{sec:reproducibility} documents why that scrutiny was necessary -- a 17-point swing on a nominally identical re-run, traced primarily to a dataset revision plus a smaller, unexplained non-determinism effect.

Given this, we resolved the crossover question directly with a dedicated campaign: retraining Standard and Fast at their exact defined configurations ($r=8$/\texttt{num\_layers=16} and $r=4$/\texttt{num\_layers=8}) across many seeds, one consistent (current) dataset vintage throughout, role \texttt{age\_5\_11}, FK$\leq$7.0 pass rate -- $n=25$/config at 1B (pre-committed, no interim look); at 3B, an initial $n=15$/config batch left the comparison marginal ($t=2.01$, $p=0.054$), so, per a plan fixed before unblinding stage two (a single top-up to $n=30$, not an open-ended stopping rule), we collected 15 further seeds/config and combine both stages below via a pre-specified equal-weight inverse-normal combination test, the valid correction for this design.

\begin{table}[!h]
  \centering
  \caption{Seed-campaign FK$\leq$7.0 pass rate (mean $\pm$ SD across independent training runs).}
  \label{tab:seedcampaign}
  \small
  \begin{tabular}{@{}lccc@{}}
    \toprule
    Config. & $n$ & Mean & SD \\
    \midrule
    Fast-1B     & 25 & 63.4\% & 6.3 \\
    Standard-1B & 25 & 52.1\% & 6.3 \\
    Fast-3B     & 30 & 52.2\% & 7.0 \\
    Standard-3B & 30 & 58.3\% & 6.9 \\
    \bottomrule
  \end{tabular}
\end{table}

On 1B, Fast beats Standard ($t=6.40$, $p=6.0\times10^{-8}$). On 3B, the naive pooled test gives $t=3.42$, $p=0.0011$, but the interim look makes this anti-conservative; combining that stage with the independent 15-seed replication collected afterward ($t=2.75$, $p=0.010$) via the combination test gives the valid $p=0.0015$ -- \textbf{the crossover survives the statistically correct test on both sides.} This campaign covers FK$\leq$7.0 pass rate only, for \texttt{age\_5\_11} only; the other four readability metrics, the \texttt{age\_12\_18} role, latency, and the classifier are single-run, as in Table~\ref{tab:res2}, though now on the same consistent dataset vintage as this campaign.

Because Standard and Fast differ systematically in response length (Table~\ref{tab:res2}), we tested whether the crossover survives controlling for length, across all 5{,}500 individual campaign responses. Fast is shorter than Standard at 1B (72 vs.\ 78 words, $p=4\times10^{-16}$) but longer at 3B (97 vs.\ 76 words, $p=9\times10^{-50}$). Residualizing FK grade on a single word-count regression pooled across configs (an equal-slopes assumption we did not test) and comparing residuals by config, \textbf{the 1B crossover survives length-adjustment} (residual gap 0.35 FK grade levels, $t=-5.95$, $p=3\times10^{-9}$), but \textbf{the 3B crossover is substantially a length effect} -- the residual gap shrinks to 0.10 FK grade levels, only marginal ($t=-1.92$, $p=0.055$). A joint logistic regression of pass/fail on config and word count agrees: the config coefficient stays large at 1B (0.35) but is near zero at 3B (0.007) once length is included. Standard's 3B advantage is better described as producing shorter, not more simply-worded, responses.

\subsection{Mechanism Sweeps: Rank and Layer Depth}\label{sec:ranksweep}

The Standard/Fast configurations differ in both rank and \texttt{num\_layers} simultaneously, so even with the crossover's existence now statistically confirmed (Section~\ref{sec:seedcampaign}), it is not yet attributable to either factor alone. The sweeps below use two seeds per configuration -- lower power than the dedicated campaign above; the reported $\pm$std is a two-point range, not a formal variance estimate. Both sweeps were retrained on the current dataset vintage (Section~\ref{sec:reproducibility}) after the original sweep was found to predate the same data revision that affected the ablation.

Fixing \texttt{num\_layers=8} and sweeping $r \in \{2,4,8,16\}$ across four architectures (Llama 1B/3B, Qwen~2 0.5B, SmolLM2~360M; two seeds; 32 runs; Table~\ref{tab:sweeps}, left) gives a more qualified picture than the original single-seed sweep suggested. \textbf{SmolLM2~360M shows a clean, monotonic rank-regime pattern} (48\%~$\to$~63\% from $r{=}2$ to $r{=}16$), consistent with depth (32 layers) driving a preference for higher rank. \textbf{Llama~1B and Qwen~0.5B, previously reported as peaking sharply at $r{=}2$, are now flat within noise} (Llama 1B: 63/59/64/62; Qwen 0.5B: 70/71/71/60), and \textbf{Llama~3B shows only a weak, noisy increase} (54\%~$\to$~62\%). We treat SmolLM2's rank sensitivity as the sweep's one robustly replicated finding, and the broader ``rank is the operative variable'' claim as suggestive, not resolved.

A companion layer sweep (fixed \texttt{rank=4}, \texttt{num\_layers}$\in\{4,8,16\}$, Llama 1B/3B, two seeds, 12 runs; Table~\ref{tab:sweeps}, right) replicates more consistently: \textbf{3B improves with depth} (57\%~$\to$~63\%, comparable to the original sweep's $+7$pt), and \textbf{1B peaks at the fewest layers} (69\% at layers=4, declining to 60\% at layers=16), matching the original sweep's conclusion despite a different exact shape. Latency scales with layer count for both sizes.

\begin{table}[!h]
  \centering
  \caption{Rank sweep (left; num\_layers=8 fixed) and layer sweep (right; rank=4 fixed) FK $\leq 7.0$ pass rate, mean $\pm$ std across two seeds (n=2; read as a range between the two runs, not a variance estimate), current dataset vintage.}
  \label{tab:sweeps}
  \scriptsize
  \setlength{\tabcolsep}{3pt}
  \begin{tabular}{@{}lcccc@{}}
    \toprule
    Model     & $r{=}2$ & $r{=}4$ & $r{=}8$ & $r{=}16$ \\
    \midrule
    Llama 1B  & 63$\pm$10 & 59$\pm$7 & \textbf{64$\pm$6} & 62$\pm$6 \\
    Llama 3B  & 54$\pm$11 & 58$\pm$3 & 57$\pm$1 & \textbf{62$\pm$11} \\
    Qwen 0.5B & 70$\pm$3 & \textbf{71$\pm$13} & 71$\pm$4 & 60$\pm$11 \\
    SmolLM2   & 48$\pm$3 & 56$\pm$6 & 62$\pm$6 & \textbf{63$\pm$4} \\
    \bottomrule
  \end{tabular}
  \qquad
  \begin{tabular}{@{}lccc@{}}
    \toprule
    Model, layers & FK$\leq$7 & Lat.(s) & Tok/s \\
    \midrule
    1B, 4  & \textbf{69\%} & 0.85 & \textbf{105.6} \\
    1B, 8  & 61\% & 0.84 & 97.4 \\
    1B, 16 & 60\% & 1.06 & 84.0 \\
    3B, 4  & 57\% & 2.73 & \textbf{45.0} \\
    3B, 8  & 60\% & 2.67 & 43.2 \\
    3B, 16 & \textbf{63\%} & \textbf{2.24} & 39.5 \\
    \bottomrule
  \end{tabular}
\end{table}

\subsection{Cross-Architecture and Capacity-Ladder Validation}\label{sec:crossarch}

To test generalization beyond Llama, we evaluate Qwen~2 0.5B and SmolLM2~360M~\cite{allal2025smollm2} (Standard/Fast, seed 42), plus the Qwen 2.5 family (0.5B, 1.5B, 3B; Standard/Fast, seed 42, 12 runs) for age~5--11, all retrained on the current dataset vintage (Table~\ref{tab:crossarch}). These results are single-seed -- far less powered than the campaign above (Section~\ref{sec:seedcampaign}) -- and should be read as exploratory, not confirmed. \textbf{SmolLM2 Standard dominates on this seed} (80\% FK, 0.87s latency), though its margin over SmolLM2 Fast (60\%) is narrower than the original sweep reported. \textbf{Qwen 2.5 shows Fast at or above Standard at every size, suggesting no crossover for that family} (56\% vs.\ 56\% at 0.5B, tied; 74\% vs.\ 66\% at 1.5B; 70\% vs.\ 50\% at 3B). Qwen 2 Fast and Qwen 2.5-0.5B Fast meet the 1.0s target with the most margin ($\sim$0.5s avg). The Classifier column reflects the original, pre-revision dataset vintage; treat these values as provisional.

\begin{table*}[!h]
  \centering
  \caption{Cross-architecture and Qwen 2.5 capacity-ladder results, age~5--11 (50 examples, seed 42, current dataset vintage; classifier is 5-fold CV, chance~=~0.50, $\dagger$original pre-revision vintage, not yet re-validated).}
  \label{tab:crossarch}
  \scriptsize
  \setlength{\tabcolsep}{3pt}
  \begin{tabular}{@{}lccccc@{}}
    \toprule
    Config. & FK$\leq$7.0 & Avg Lat. & $<$1.0s & Tok/s or Tok & Classifier$^\dagger$ \\
    \midrule
    Qwen 2 Std.       & 68\% & 0.54\,s & 100\% & 83 tok  & $0.940 \pm 0.049$ \\
    Qwen 2 Fast       & 62\% & 0.51\,s & 96\%  & 94 tok  & $0.960 \pm 0.058$ \\
    SmolLM2 Std.      & 80\% & 0.87\,s & 72\%  & 82 tok  & $0.950 \pm 0.032$ \\
    SmolLM2 Fast      & 60\% & 1.30\,s & 56\%  & 137 tok & $0.920 \pm 0.040$ \\
    Qwen2.5-0.5B Fast & 56\% & 0.45\,s & 100\% & 176.4   & $0.950 \pm 0.032$ \\
    Qwen2.5-0.5B Std. & 56\% & 0.57\,s & 100\% & 152.1   & $0.940 \pm 0.097$ \\
    Qwen2.5-1.5B Fast & 74\% & 1.07\,s & 58\%  & 76.1    & $\mathbf{0.980 \pm 0.024}$ \\
    Qwen2.5-1.5B Std. & 66\% & 1.28\,s & 6\%   & 67.8    & $0.890 \pm 0.073$ \\
    Qwen2.5-3B Fast   & 70\% & 1.73\,s & 4\%   & 43.0    & $0.970 \pm 0.040$ \\
    Qwen2.5-3B Std.   & 50\% & 2.18\,s & 0\%   & 39.9    & $0.930 \pm 0.068$ \\
    \bottomrule
  \end{tabular}
\end{table*}

\section{Methods}\label{sec:methodology}

\subsection{System Architecture}\label{sec:architecture}

Raw text first passes through \textit{guard-bear}, a binary input classifier; unsafe input is blocked before reaching the response model. Safe input is routed by an age-gate, supplied explicitly by the calling VR application, to one of two LoRA adapters on a shared frozen 4-bit quantized base model. Neither the adapters nor the guard classifier is fused into its base model. Emotional-support queries are in scope only as reassurance/redirection; clinical assessment and crisis intervention are out of scope. No system prompt is used, so the adapter weights encode register directly.

\textit{guard-bear} classifies raw text as \textbf{safe} (a legitimate pediatric query, including benign out-of-scope curiosity) or \textbf{unsafe} (prompt injection, jailbreak attempts, adult/parental queries, gibberish, inappropriate content, or social engineering), and is a full parameter fine-tune of a small (86M) classifier, unlike the response model. The design target is high recall: a missed unsafe input is materially worse than an incorrectly blocked legitimate query, motivating the $\geq 0.97$ recall target.

\subsection{Datasets}\label{subsec:responsedataset}

The response-model dataset was generated by Claude Opus 4.6 in extended thinking mode from the Children's Hospital Patient and Family Guide~\cite{hospitalguide}, across five categories (Emotional Reassurance, FAQs/General Curiosity, Hospital Rules and Routines, What to Expect, Who Are These People?), each age-stratified (5--11, 12--18). An initial 1000-example release (200/category, no \texttt{edge\_cases}) was superseded by a quality-curation pass adding a training-only \texttt{edge\_cases} category (50) and 73 diversified examples, for 1123 total, split 562/561 (Section~\ref{sec:trainprotocol}), and replacing the hand-curated validation set with 100 new examples; Section~\ref{sec:reproducibility} documents how this revision affected earlier results. Examples are instruction/response pairs tagged by role and category, without a system prompt, so the model encodes register in its weights.

\textit{guard-bear}'s dataset combines synthetic safe/unsafe examples with unsafe examples pulled and filtered from public prompt-injection and jailbreak datasets: 4{,}703 examples (2{,}404 unsafe, 2{,}299 safe) across 14 subcategories -- seven safe (e.g.\ \texttt{hospital\_information}, \texttt{emotional\_support}, \texttt{pre\_procedure\_anxiety}) and seven unsafe (e.g.\ \texttt{prompt\_injection}, \texttt{adult\_parental\_query}, \texttt{gibberish\_malformed}); full taxonomy in the repository. ``Safe'' means \textit{not actively harmful or adversarial}, not ``in the response model's scope''; benign out-of-scope questions are handled downstream by the response model itself. The corpus is split 3{,}476/614/613 for training, validation, and held-out test.

\subsection{Base Models, LoRA Configuration, and Training Protocol}\label{sec:trainprotocol}

We use the Llama 3.2 Instruct family~\cite{dubey2024llama3} at 1B/3B, 4-bit quantized via \texttt{mlx\_lm}~\cite{apple2023mlx} on Apple Silicon; full-parameter fine-tuning of the 3B model ($\sim$30GB bfloat16) exceeds the 16GB ceiling by nearly 2$\times$, making LoRA the only viable strategy. \textit{guard-bear} fine-tunes all 86M parameters via the Metal Performance Shaders (MPS) backend, well within the same budget.

LoRA~\cite{hu2022lora} freezes a pre-trained weight matrix $W_0 \in \mathbb{R}^{d\times k}$ and expresses updates as $\Delta W = BA$, $B \in \mathbb{R}^{d\times r}$, $A \in \mathbb{R}^{r\times k}$, $r \ll \min(d,k)$, applied to all linear projections in the targeted layers with \texttt{scale=20.0} held constant (a raw output multiplier in \texttt{mlx\_lm}, not an alpha/rank ratio). \textbf{Standard} ($r=8$, \texttt{num\_layers=16}) covers the upper 16 layers (57\% of 3B, 100\% of 1B); \textbf{Fast} ($r=4$, \texttt{num\_layers=8}) covers the upper 8 (29\%/50\%). The two configurations differ in both rank and layer count simultaneously; Sections~\ref{sec:seedcampaign}--\ref{sec:ranksweep} confirm the resulting crossover is real and partially decompose its mechanism.

Response-model adapters train with Adam, cosine-decay LR (peak $1\times10^{-5}$, 100-step warmup, 600 total steps), batch size 4, max sequence length 768, LoRA dropout 0.0, loss masked to assistant tokens only; all eight runs use seed 42. \textit{guard-bear} fine-tunes \texttt{meta-llama/Prompt-Guard-86M}~\cite{promptguard2024} (DeBERTa-v2-based~\cite{he2021deberta}) on the full training split, seed 42, with the classification threshold tuned post-hoc on the validation split to the lowest value meeting the $\geq 0.97$ recall target, selecting \textbf{0.08}; the test split is used only for final evaluation.

\noindent\textbf{Reproducibility and dataset vintage.}\label{sec:reproducibility} Several results here required a dedicated reproducibility investigation before they could be trusted. During rank-sweep construction, a routine re-run of the nominal Fast-1B configuration (identical model, data, seed) produced an FK$\leq$7.0 pass rate 17 points lower than the original single run reported in an earlier draft. A confounding \texttt{mlx}/\texttt{mlx\_lm} library-version change was ruled out; the training and validation data had also been revised between runs (Section~\ref{subsec:responsedataset}), and re-scoring the original adapter on the revised validation set alone closed most of the gap. A smaller, genuine non-determinism effect (16--18 points, data held fixed) remains and is not fully explained. A closer audit (timestamps against data modification times) then found the rank sweep, layer sweep, and initial cross-architecture results had also, unexpectedly, predated the revision; all were retrained on the current vintage before being reported here (Sections~\ref{sec:seedcampaign}--\ref{sec:crossarch}), except Table~\ref{tab:crossarch}'s Classifier column, still provisional. This episode directly motivated the dedicated multi-seed campaign in Section~\ref{sec:seedcampaign}: a single-run, single-vintage result on this pipeline cannot be trusted without independent scrutiny, a standard applied throughout the results above.

\subsection{Evaluation Framework}

For the response model: \textbf{readability} via \texttt{textstat}~\cite{bansal2023textstat} (Flesch-Kincaid~\cite{kincaid1975fk} primary, pass/fail at FK~$\leq$~7.0 for 5--11 outputs; SMOG~\cite{mclaughlin1969smog}, Gunning Fog~\cite{gunning1952fog}, Coleman-Liau~\cite{coleman1975cli}, type-token ratio); \textbf{latency/throughput} under sequential, single-request conditions, 1.0s average threshold; \textbf{inter-role style separation} via TF-IDF+logistic-regression (scikit-learn~\cite{pedregosa2011sklearn}, 5-fold CV, accuracy~$\geq$~0.90, chance~=~0.50). These are proxies, not direct measures, of age-adapted communication: readability formulas capture sentence- and word-length complexity, not conceptual difficulty or emotional register, and the style classifier captures lexical distinguishability, not developmental appropriateness (Section~\ref{sec:discussion}). For \textit{guard-bear}: unsafe recall (primary), precision, F1, accuracy, and ROC-AUC, plus confusion matrix and per-subcategory recall, against the unmodified base model.

\textbf{Ethical approval declarations.} Not applicable: no experiments on live vertebrates, humans, or human tissue; all data are synthetic or drawn from public reference/prompt-injection materials, with no patient contact or identifying information.

\section{Discussion}\label{sec:discussion}

For \textit{guard-bear}, the central finding is that \textbf{a general-purpose safety classifier cannot be assumed to transfer to a new deployment population without evaluation}; the gap here is severe, worse-than-random discrimination, not merely suboptimal thresholding. This matters directly for deployment: the response-model results below assume in-scope, benign input, and \textit{guard-bear} is what makes that assumption safe here. An age-adaptive response model without an input safety gate is not, by itself, deployable, and Section~\ref{sec:deployment} shows both are practical together on consumer hardware.

\textbf{Fast-1B is the best Llama-family configuration} for VR deployment: the only Llama option meeting the 1.0s latency target with strong readability (FK 6.5 avg) and style separation (0.960); SmolLM2 Standard (Section~\ref{sec:crossarch}) beats it on FK pass rate (80\% vs.\ 62\%) and is a strong alternative where architecture is unconstrained. Beyond this deployment choice, the response model's most interesting outcome was not planned: the \textbf{crossover} discovered during development and statistically confirmed (Section~\ref{sec:seedcampaign}) suggests small-model LoRA configuration below 7B does not follow rank-selection guidance developed at larger scales -- though length-adjustment shows only the 1B side is length-independent; the 3B side is substantially Standard producing shorter, not more simply-worded, responses. Its mechanism is likewise only partially resolved: the rank sweep's capacity-regularization story replicates cleanly for SmolLM2 but not Llama or Qwen at the sample sizes tested (Section~\ref{sec:ranksweep}). These gains are necessary but not sufficient evidence of genuinely age-adapted communication (Section~\ref{sec:methodology}): a response can be short and lexically distinct from the other age group's outputs while still conceptually or emotionally mismatched to its audience, a gap only human or clinical evaluation can close.

\section{Conclusion}\label{sec:conclusion}

This paper shows two things. First, a general-purpose input-safety classifier cannot be assumed to transfer to a new, vulnerable population: \texttt{Prompt-Guard-86M}'s unmodified performance on pediatric-context queries is worse than random, resolved only by dedicated fine-tuning to near-perfect discrimination (ROC-AUC 0.999) -- a finding with implications beyond this deployment. Second, a complete age-adaptive response and safety pipeline for pediatric hospital communication is achievable on one consumer machine (Apple Mac Mini M4, 16~GB), producing measurable improvements in automated readability and style-separation proxies over zero-shot base models (Abstract).

During development of the response model, we also uncovered and rigorously validated an unexpected LoRA configuration-ordering crossover between 1B and 3B Llama models (Section~\ref{sec:seedcampaign}) after an initial single-seed observation proved non-reproducible at face value. A rank sweep and layer sweep sought to attribute this to rank (capacity regularization) and layer depth; the layer-depth contribution replicates reasonably well on independent retraining, but the rank-regime mechanism replicates cleanly only for SmolLM2 and is markedly weaker for Llama and Qwen at two-seed power -- the crossover's existence is established, its precise cause only partially so. Neither the response model nor \textit{guard-bear} alone is deployable: the response model presumes benign input, and \textit{guard-bear} is what makes that presumption safe.

Limitations: Section~\ref{sec:reproducibility} documents this pipeline's non-determinism and dataset-vintage confound, accounted for above; one piece remains provisional (Table~\ref{tab:crossarch}'s Classifier column). The rank sweep's mechanism claim is qualified (Section~\ref{sec:ranksweep}): SmolLM2's rank sensitivity replicates cleanly, Llama and Qwen do not, at two-seed power -- plausible, not resolved. Length-adjustment was untested on the sweeps' architectures; single-seed perplexity does not track the crossover; models above 3B remain untested. More fundamentally, training and validation data share one generative process, and every response-model metric is fully automated: no human rater or clinician has judged any output. This closed evaluation loop confirms internal, not external, validity -- human evaluation is necessary before any deployment claim. The base-model comparison is zero-shot with no system prompt, the weakest baseline; neither component has clinical validation with pediatric patients or staff. Reported latency excludes \textit{guard-bear}'s pass; end-to-end latency is unmeasured, nor has \textit{guard-bear}'s dataset been red-teamed or checked for train/test leakage.

\backmatter

\bmhead{Supplementary information}

Full validation-set outputs, dataset files, and training/evaluation scripts are available at \url{https://github.com/abarros6/small-bear}; no separate supplementary bundle accompanies this submission.

\bmhead{Acknowledgements}

Not applicable.

\section*{Declarations}

\begin{itemize}
\item \textbf{Funding:} No funding to disclose.
\item \textbf{Conflict of interest/Competing interests:} The authors declare no conflicts of interest.
\item \textbf{Ethics approval and consent to participate; consent for publication:} Not applicable; no human subjects, patient data, or clinical trial. All data are synthetic or from public reference/prompt-injection materials.
\item \textbf{Data, materials, and code availability:} Datasets, model configurations, and all training/inference/evaluation code are available at \url{https://github.com/abarros6/small-bear}.
\item \textbf{Author contribution:} S.d.R.\ and R.E.\ supervised the research. A.B.\ ran all experiments, designed \textit{guard-bear}, and wrote the paper except the Introduction, written by K.K. All authors reviewed and approved the final manuscript.
\end{itemize}

\begin{thebibliography}{99}

\bibitem{zhou2023lima}
C.~Zhou, P.~Liu, P.~Xu, S.~Iyer, J.~Sun, R.~Maddela, X.~Peng, S.~Shrivastava, M.~Lewis, L.~Zettlemoyer, and O.~Levy, ``LIMA: Less Is More for Alignment,'' in \textit{Advances in Neural Information Processing Systems (NeurIPS)}, 2023.

\bibitem{dubey2024llama3}
A.~Dubey et al., ``The Llama 3 Herd of Models,'' \textit{arXiv preprint arXiv:2407.21783}, Meta AI, 2024.

\bibitem{apple2023mlx}
Apple, ``MLX: Efficient and Flexible Machine Learning on Apple Silicon,'' 2023. [Online]. Available: \url{https://github.com/ml-explore/mlx}

\bibitem{hu2022lora}
E.~J.~Hu, Y.~Shen, P.~Wallis, Z.~Allen-Zhu, Y.~Li, S.~Wang, L.~Wang, and W.~Chen, ``LoRA: Low-Rank Adaptation of Large Language Models,'' in \textit{Proc. Int. Conf. Learning Representations (ICLR)}, 2022.

\bibitem{kincaid1975fk}
J.~P.~Kincaid, R.~P.~Fishburne, R.~L.~Rogers, and B.~S.~Chissom, ``Derivation of New Readability Formulas for Navy Enlisted Personnel,'' Research Branch Report 8-75, Naval Technical Training Command, 1975.

\bibitem{mclaughlin1969smog}
G.~H.~McLaughlin, ``SMOG Grading --- A New Readability Formula,'' \textit{J. Reading}, vol.~12, no.~8, pp.~639--646, 1969.

\bibitem{gunning1952fog}
R.~Gunning, \textit{The Technique of Clear Writing}. New York, NY, USA: McGraw-Hill, 1952.

\bibitem{coleman1975cli}
M.~Coleman and T.~L.~Liau, ``A Computer Readability Formula Designed for Machine Scoring,'' \textit{J. Applied Psychology}, vol.~60, no.~2, pp.~283--284, 1975.

\bibitem{bansal2023textstat}
S.~Bansal, ``textstat: Python Library to Calculate Readability Statistics of a Text Object'' (v0.7.13), 2023. [Online]. Available: \url{https://pypi.org/project/textstat/}

\bibitem{dettmers2023qlora}
T.~Dettmers, A.~Pagnoni, A.~Holtzman, and L.~Zettlemoyer, ``QLoRA: Efficient Finetuning of Quantized LLMs,'' in \textit{Advances in Neural Information Processing Systems (NeurIPS)}, 2023.

\bibitem{biderman2024lora}
D.~Biderman, J.~G.~Ortiz, J.~Portes, M.~Paul, P.~Greengard, C.~Jennings, D.~King, S.~Havens, V.~Chiley, J.~Frankle, C.~Blakeney, and J.~P.~Cunningham, ``LoRA Learns Less and Forgets Less,'' \textit{Transactions on Machine Learning Research (TMLR)}, 2024.

\bibitem{pedregosa2011sklearn}
F.~Pedregosa et al., ``Scikit-learn: Machine Learning in Python,'' \textit{J. Machine Learning Research}, vol.~12, pp.~2825--2830, 2011.

\bibitem{allal2025smollm2}
L.~B. Allal et al., ``SmolLM2: When Smol Goes Big,'' \textit{arXiv preprint arXiv:2502.05737}, Hugging Face, 2025.

\bibitem{hospitalguide}
``Children's Hospital Patient and Family Guide,'' Children's Hospital London Health Sciences Centre. \url{https://www.lhsc.on.ca/childrens-hospital/childrens-hospital-patient-and-family-guide} (accessed Mar.\ 26, 2026).

\bibitem{promptguard2024}
Meta AI, ``Prompt-Guard-86M,'' 2024. [Online]. Available: \url{https://huggingface.co/meta-llama/Prompt-Guard-86M}

\bibitem{he2021deberta}
P.~He, X.~Liu, J.~Gao, and W.~Chen, ``DeBERTa: Decoding-enhanced BERT with Disentangled Attention,'' in \textit{Proc. Int. Conf. Learning Representations (ICLR)}, 2021.

\bibitem{rao2025pediatric}
A.~S.~Rao, A.~Mazumder, E.~Roux, C.~Young, E.~Bott, J.~Wang, M.~Kochis, A.~Stetson, A.~Butler, S.~Hilker, and M.~D.~Succi, ``Designing Patient-Centered Communication Aids in Pediatric Surgery Using Large Language Models,'' \textit{Journal of Pediatric Surgery}, vol.~60, no.~12, art.~162654, 2025.

\bibitem{hassan2025childconv}
S.~Z.~Hassan, P.~Halvorsen, M.~S.~Johnson, and P.~Lison, ``Evaluating LLMs on Generating Age-Appropriate Child-Like Conversations,'' \textit{arXiv preprint arXiv:2510.24250}, 2025.

\bibitem{nayeem2024kidlm}
M.~T.~Nayeem and D.~Rafiei, ``KidLM: Advancing Language Models for Children -- Early Insights and Future Directions,'' in \textit{Proc. 2024 Conf. Empirical Methods in Natural Language Processing (EMNLP)}, pp.~4813--4836, Miami, FL, USA, 2024.

\bibitem{ryu2021vr}
J.-H.~Ryu, J.-W.~Park, S.~I.~Choi, J.~Y.~Kim, H.~Lee, H.-J.~Yoo, and S.-H.~Han, ``Virtual Reality vs. Tablet Video as an Experiential Education Platform for Pediatric Patients Undergoing Chest Radiography: A Randomized Clinical Trial,'' \textit{J. Clin. Med.}, vol.~10, no.~11, art.~2486, 2021.

\bibitem{kidrails2025}
Arcee AI, ``KidRails: Replicating a Child-Safe LLM Model,'' 2025. [Online]. Available: \url{https://github.com/arcee-ai/KidRails}

\end{thebibliography}

\end{document}
